\justify

\section{Literature Survey}
\justify
The following table presents a comparative survey of key research works in the areas of
AI-based plant identification, large language models in healthcare, digital Ayurveda,
and cross-platform mobile development---the foundational domains of AyurSpace.

\begin{table}[H]
\caption{Literature Survey}
\label{tab:lit_survey}
{\small\setlength{\tabcolsep}{4pt}\renewcommand{\arraystretch}{0.95}
\centering
\begin{tabular}{|>{\centering\arraybackslash}p{0.7cm}|
                 >{\centering\arraybackslash}p{3.3cm}|
                 >{\centering\arraybackslash}p{0.8cm}|
                 >{\centering\arraybackslash}p{2.5cm}|
                 >{\centering\arraybackslash}p{3.3cm}|
                 >{\centering\arraybackslash}p{3.3cm}|}
\hline
\textbf{Sr.~No.} & \textbf{Title of Paper} & \textbf{Year} &
\textbf{Author} & \textbf{Advantages} & \textbf{Limitations} \\
\hline
1 &
Automated Identification of Medicinal Plants Using Deep CNNs~[11] &
2023 & Kumar et al. &
High accuracy in species classification; effective on high-resolution images &
Limited training data for rare Indian medicinal plants; poor performance in varied lighting \\
\hline
2 &
PlantNet: A Citizen Science Approach to Plant Identification~[12] &
2021 & Bonnet et al. &
Large crowd-sourced global dataset; supports multi-organ recognition &
European flora bias; lacks Ayurvedic or medicinal use context \\
\hline
3 &
Large Language Models Encode Clinical Knowledge~[4] &
2023 & Singhal et al. &
Achieves expert-level medical Q\&A; generalises across medical specialties &
Western medicine bias; lacks alignment with traditional Indian medicine \\
\hline
4 &
Conversational Agents in Healthcare: A Systematic Review~[13] &
2022 & Laranjo et al. &
Enhances patient engagement; supports remote wellness monitoring &
Prone to context drift; limited cultural and language sensitivity \\
\hline
5 &
A Digital Framework for Tridosha-Based Health Assessment~[14] &
2021 & Patel \& Sharma &
Formalises Prakriti classification for digital systems; reproducible scoring &
Subjective question responses; no universally accepted digital scale \\
\hline
6 &
AI-Powered Personalized Wellness Applications: A Review~[15] &
2022 & Mehta et al. &
Personalised recommendations shown to improve user health outcomes &
Significant data privacy concerns; risk of over-reliance on AI suggestions \\
\hline
\end{tabular}}
\end{table}

\newpage

\section{Related Work}

\noindent\textbf{Computer Vision in Botany:}\\
The area of automated plant identification has grown significantly with the advent of
Convolutional Neural Networks (CNNs). Initial endeavors employed leaf-shape descriptors
and edge detection algorithms~[5]. Modern methods use Deep Learning
architectures such as ResNet-50 and MobileNetV3. \textbf{Pl@ntNet} excels with European
plants but struggles with Indian medicinal herbs. \textbf{LeafSnap} focuses on leaf
curvature for tree species. These systems operate in a ``Botanical Silo''---they provide
only a Latin binomial, whereas for Ayurveda, identification is just the beginning of the
consultation.

\vspace{5mm}
\noindent\textbf{LLMs in Healthcare:}\\
Generative AI has shown promise in combining medical knowledge. Singhal et al.~[4]
showed that Med-PaLM could pass the USMLE. However, applying general-purpose LLMs to
traditional knowledge bases frequently leads to ``alignment drift,'' prioritizing Western
medical interpretations over traditional logic. There is no dedicated ``Ayur-PaLM,''
making prompt engineering essential to constrain LLMs like Gemini to act as reliable
Ayurvedic experts.

\vspace{5mm}
\noindent\textbf{Digital Approaches to Ayurveda:}\\
Previous research in ``Computational Ayurveda'' has focused on structural diagnostics
and static databases---Pulse Diagnosis (\textit{Nadi Pariksha}) using piezoelectric
sensors, and Tongue Diagnosis (\textit{Jivha Pariksha}) using image processing. Existing
Ayur Apps are simple CRUD applications without active AI context or personalized
adaptability.

\vspace{5mm}
\begin{table}[H]
\caption{AyurSpace Compared to Other Solutions}
\begin{center}
\begin{tabular}{|c|c|c|c|c|}
\hline
\textbf{Feature} & \textbf{PlantNet} & \textbf{Google Lens} & \textbf{Ayur Apps} & \textbf{AyurSpace} \\
\hline
\textbf{Visual ID} & High & High & None & \textbf{High} \\
\hline
\textbf{Context} & None & Limited & Static & \textbf{Dynamic} \\
\hline
\textbf{Dosha Logic} & None & None & Basic & \textbf{Adaptive} \\
\hline
\textbf{Interaction} & Passive & Passive & Passive & \textbf{Chat} \\
\hline
\textbf{Safety} & N/A & Low & Medium & \textbf{High} \\
\hline
\end{tabular}
\end{center}
\end{table}

\newpage
\section{Existing Systems}
\justify
In the existing landscape of plant identification and Ayurvedic wellness applications,
monitoring and guidance are primarily performed using fragmented, single-purpose tools
that lack the integration necessary for safe and personalized Ayurvedic practice. Users
are required to switch between separate apps for plant identification, health tracking,
and dietary guidance, with no unified platform providing end-to-end Ayurvedic reasoning.

\textbf{Pl@ntNet} and \textbf{iNaturalist} are leading plant identification apps driven
by crowd-sourced datasets and CNNs. While accurate for European flora, they operate in
a ``Botanical Silo''---providing only a Latin binomial with no medicinal context,
making them fundamentally unsuitable for Ayurvedic practice.

\textbf{Google Lens} offers general-purpose visual recognition but does not specialize
in Indian medicinal herbs and frequently misidentifies visually similar species. Its
information is limited to Wikipedia-level data with no Ayurvedic property output or
dosha analysis.

\textbf{Existing Ayurvedic mobile applications} such as Jiva Ayurveda and NirogStreet
are essentially static databases of herbs. They lack AI-based plant scanning, real-time
LLM reasoning, and adaptive dosha scoring, offering no conversational interface for
personalized guidance.

\textbf{General-purpose LLM chatbots} (ChatGPT, Gemini standalone) are prone to
hallucinating plant properties and fabricating citations~[4]. Without a grounding
visual identification step, they cannot be trusted for plant-specific medical guidance.

These limitations collectively highlight the need for an integrated, AI-powered mobile
system that combines accurate visual plant identification, structured Ayurvedic property
retrieval, personalized \textit{Prakriti} assessment, and a safety-constrained
conversational interface---the exact gap that AyurSpace addresses.

\newpage
\section{Problem Statement}
\justify
The digitization of Ayurveda introduces specific computational challenges that current
solutions do not adequately resolve:
\begin{enumerate}[label={\textbf{\arabic*.}}, leftmargin=*]
    \item \textbf{Taxonomic Ambiguity}: Numerous distinct species possess identical
    nomenclature. For example, ``Brahmi'' can mean either \textit{Bacopa monnieri} or
    \textit{Centella asiatica}, depending on region. A text-only search is insufficient
    and dangerous.

    \item \textbf{Visual Similarity}: Many beneficial plants resemble harmful ones.
    For instance, the Solanaceae family contains both poisonous plants and edible
    vegetables, requiring accurate computer vision to distinguish them.

    \item \textbf{Contextual Disconnection}: Current plant identification apps (like
    PlantNet) provide Linnaean taxonomy but no Ayurvedic context. Knowing a plant is
    \textit{Tinospora cordifolia} is not useful without knowing its \textit{Guduchi}
    properties and dosage.

    \item \textbf{Generative Hallucinations}: LLMs like GPT-4 can provide Ayurvedic
    advice but are prone to hallucinating citations and properties. Direct visual LLM
    identification is presently unreliable for essential medical safety~[4].
\end{enumerate}

To address these challenges, we present \textbf{AyurSpace}. Our contributions are:
\begin{itemize}
    \item \textbf{Hybrid Neuro-Symbolic Architecture}: A pipeline that separates
    ``Identification'' (specialized CNNs) from ``Reasoning'' (Semantic LLMs), minimizing
    hallucination risks.
    \item \textbf{Digitized Dravyaguna Ontology}: Medicinal plant properties in a
    searchable, object-oriented schema with algorithmic safety checks.
    \item \textbf{Quantitative Tridosha Assessment}: \textit{Prakriti Pariksha}
    formalized into a discrete-math replicable scoring algorithm.
    \item \textbf{Open Source Mobility}: Cross-platform Flutter implementation for
    accessibility on affordable devices common in developing countries.
\end{itemize}

